% !TEX root = ../main.tex
\section{Grounding, Context, and Applicability}
\label{sec:grounding-context}

Definition~\ref{def:grounded-evidence} listed seven conditions on evidence without
explaining them. This section takes them in turn and identifies the generalisation error
that makes applicability the hardest of the seven to enforce. The resolved context $c$
that four of them depend on was defined in Section~\ref{sec:resolved-context}.

\subsection{Illicit universalization}
\label{sec:illicit}

Relation \eqref{eq:no-transfer} names a failure that deserves its own label,
because treating it as ordinary inaccuracy leads to the wrong remedy.

\emph{Illicit universalization} is the treatment of a claim that holds under one
jurisdiction, time, population, policy, or use as though it held generally, because
the relevant context was never resolved. The claim is not fabricated. It is not
even, in its own frame, wrong. It has been detached from the frame that made it
assessable, and then asserted into a frame where its truth value is simply a
different question.

This is worth insisting on because it reframes a large fraction of what gets
reported as hallucination. Many apparent failures are not invented facts. They are
unresolved-reference failures, temporal-validity failures, jurisdiction failures,
or applicability failures: information that is correct somewhere, stated without a
valid binding to the context in which reliance is proposed. The diagnostic question
is not ``is this true?'' but ``true \emph{where}, \emph{when}, \emph{for whom}, and
\emph{for what use}?''---and a system that cannot represent the question cannot be
said to have answered it.

The data-architecture analogue is exact, and worth stating because it makes the
failure recognisable rather than novel. A measure without its dimension keys is not
a fact; it is a number. Revenue is meaningless without period, entity, and currency,
and no data architect would accept it into a warehouse unqualified. Illicit
universalization is the same defect at the level of an assertion: a claim detached
from the dimensions that make it assessable, presented as though those dimensions
were universal rather than merely unstated. What ordinary dimensional modelling
enforces by schema, these systems do not enforce at all, because by
\eqref{eq:signature} there is no schema to enforce it with.

There is a structural reason this failure is native to the architecture rather than
incidental. By \eqref{eq:signature} the system's input is a token sequence, and a
token sequence carries no context index. A question that omits its jurisdiction is
not, to the system, an incomplete question; it is a complete string. Nothing in
\eqref{eq:objective} rewards noticing that a binding is missing, and by
Section~\ref{sec:rlhf} preference optimisation actively rewards answering anyway.
The system is not failing to resolve context. It has no representation of context
to resolve, and no mechanism that could register the omission.

\subsection{The seven dimensions}
\label{sec:six-dimensions}

\paragraph{Authentic.} The source is what it purports to be. In practice this means
retrieval from a controlled corpus whose contents have known origin, not from a
scrape whose provenance is a URL. Authenticity is a property an organisation
establishes by ingestion controls, and it is the dimension most familiar from
ordinary data governance.

\paragraph{Versioned.} The source version, its effective date, and the corpus state
$t_k$ at retrieval can all be identified. A regulation without an effective date is
not evidence about what was required at a decision date; it is evidence about what
some document said at some point. Both timestamps are needed, for different purposes:
the effective date establishes validity at $t_v$, and $t_k$ establishes what the
system could have seen. This is the dimension that permits replay, and without it the record itself cannot
support reconstruction: whatever an investigator eventually pieces together from
archives and change logs, it will not be a reconstruction from the system's own
evidence.

\paragraph{Authoritative.} The source has the standing required for the decision
type. A vendor summary of a regulation and the regulation itself may agree on
content and differ entirely in standing, and which one is required is a function of
the decision, not of the content. Authority is a relation between a source and a
decision class, not a scalar attached to a document---which is why an authority
ranking that is global rather than use-conditional will misclassify.

\paragraph{Applicable.} The source governs or materially bears on $c$. Read
component-wise, this requires the source to govern in jurisdiction $j$, to be in force
at valid time $t_v$---interval containment, per
Section~\ref{sec:resolved-context}, so neither not-yet-enacted nor
already-superseded---to bear on the material facts $f$, and to cover the intended use
$u$. This is where \eqref{eq:no-transfer} bites, and it is the dimension standard
retrieval most reliably fails, for reasons developed in
Section~\ref{sec:retrieval-insufficient}.

\paragraph{Supports.} The source establishes $p$ at the strength asserted.
Section~\ref{sec:entails-supports} distinguishes the grades of support this covers.

\paragraph{Sufficient.} The source addresses every condition the claim requires,
not merely some of them. A rule imposing five conditions is not satisfied by evidence
speaking to three, and an assertion resting on partial evidence is grounded in every
other dimension while being inadequate in the one that matters. Sufficiency and
support are separated because their remedies differ: insufficient support is fixed by
retrieving \emph{more}, inadequate support by retrieving \emph{better}.

\paragraph{Traceable.} A reviewer can reconstruct the evidence-to-assertion
relation: which passage, which version, retrieved at which corpus state, supporting
which claim. Traceability is a property of the record rather than of the evidence, and
it is what turns the other six into an auditable finding rather than an assertion
about the system's internals.

Note that four of the seven---$\Authoritative$, $\Applicable$, $\Supports$, and
$\Sufficient$---are relations to $c$ rather than properties of the document. A corpus
can be perfectly authentic, fully versioned, and completely traceable while failing
all four, and a governance programme that treats grounding as a corpus property will
build exactly that system. The no-defeater condition of \eqref{eq:nodefeater} is
stronger still: it is a relation to the \emph{entire base}, so no examination of a
single retrieved source can establish it.

\subsection{Entailment and support}
\label{sec:entails-supports}

$\Supports(E, p, c)$ covers a range, and collapsing it is a common source of
overclaiming. At the strong end,
\begin{equation}
  \Entails(E, p, c) \;\Rightarrow\; \Supports(E, p, c),
  \label{eq:entails}
\end{equation}
where $E$ together with the governing rules of $c$ settles $p$. A capital
calculation that follows deterministically from a rule and a set of positions is of
this kind, and so is a policy question whose answer is stated in the policy.

Most enterprise questions are weaker. Evidence may make $p$ substantially more
likely without settling it; it may support a qualified version of $p$ but not the
unqualified assertion; it may support $p$ for a class of cases of which $c$ is
plausibly but not demonstrably a member. None of this is a defect. It is the normal
condition of decision-making under real evidence.

The requirement is not that every assertion be entailed. It is that the kind and
strength of support be \emph{explicit and proportional to the consequence of
reliance}. A system that reports ``supported'' without distinguishing entailment
from suggestive correlation has not given a reviewer what they need, and has
created the specific failure in which an assertion adequate for one decision class
is relied upon for another. This is cell (iii) in its most defensible form: the
evidence was authentic, applicable, and authoritative, and it was adequate for a
weaker claim than the one asserted.

\subsection{Why retrieval and citation are not enough}
\label{sec:retrieval-insufficient}

The anticipated objection to this entire section is that retrieval-augmented
generation \citep{lewis2020} already addresses it. Retrieval does help, materially,
and no argument here counts against using it. It does not establish grounding in
the sense of Definition~\ref{def:grounded-evidence}, and the gap is specific.

Standard retrieval ranks passages by semantic similarity to the query and places
the top results in the context window. Consider what that procedure establishes
against the six dimensions. It can establish authenticity, if the corpus is
controlled. It can establish traceability, if the retrieval event is logged. It
establishes nothing about versioning unless version metadata is carried and used,
nothing about authority unless authority is represented and conditioned on the
decision type, and---critically---nothing about applicability, because semantic
similarity to the query is not applicability to the context.

That last point deserves to be stated as a slogan, because it is the practical core
of this paper:

\begin{quote}
Retrieval is not applicability. Citation is not entailment. And authoritative
evidence is not automatically governing evidence for the case at hand.
\end{quote}

The nearest existing measurement is worth naming, since retrieval evaluation does
score something adjacent. Context-relevance metrics assess whether retrieved evidence
is relevant \emph{to the query}. Applicability asks whether it \emph{governs the
case}. A paediatric question that retrieves adult guidance scores as highly
relevant---the topic matches exactly---and the guidance does not govern. Relevance to a
query and authority over a decision are different relations, and only the first is
measured.

Each clause of the slogan names a distinct substitution that systems make routinely. A
passage
about paediatric dosing and a passage about adult dosing are near-neighbours in
embedding space; that is what makes similarity search work, and it is exactly why
similarity search will surface the wrong one. A citation demonstrates that a
document was consulted, not that it supports the claim attached to it, and
citation-to-claim alignment is a separate check that standard pipelines do not
perform. And a source may be the most authoritative document in the corpus and
still not govern this decision.

Nor does $\Sensitive$ come for free. Placing evidence in the context raises the
probability of consistent continuations; it does not make the assertion conditional
on the evidence having supported it. By Section~\ref{sec:missing} retrieved text is
more tokens, not a precondition, and a system may retrieve an authoritative passage
and then answer from parametric memory, agreeing with it or not.

The conclusion is not that retrieval fails. It is that retrieval supplies one of
six dimensions reliably, two more if metadata is carried and logged, and none of
the three that are relations to $c$. The remaining work is architecture, and
Section~\ref{sec:governed-rag} sets out what it consists of.
